% TOP_PROBLEM: {"id": 1142, "title": "The Painlevé Conjecture", "category_id": 11, "category": {"id": 11, "name": "dynamical_systems", "display_name": "Dynamical Systems", "description": "Problems about long-term behavior of deterministic systems, Hamiltonian dynamics, and periodic orbits.", "slug": "dynamical-systems", "order_index": 11, "created_at": "2026-07-31T15:26:25.671Z"}, "status": "open"}
% FIRST_POSTED: null
% ENTRY_KIND: statement_only
% Imported from: batch04.tex; UnsolvedMath ID 1142
\documentclass[11pt]{article}
\usepackage[margin=25mm]{geometry}
\usepackage{fontspec}
\setmainfont{Latin Modern Roman}
\usepackage{amsmath,amssymb,mathtools}
\usepackage{unicode-math}
\setmathfont{Latin Modern Math}
\usepackage{xurl,hyperref}
\hypersetup{colorlinks=true,urlcolor=blue}
\setcounter{secnumdepth}{0}
\setlength{\parindent}{0pt}
\setlength{\parskip}{5pt}
\setlength{\emergencystretch}{3em}
\newcommand{\sref}[2]{\hyperlink{source:#1:#2}{[S#2]}}
\newcommand{\cataloguescope}[1]{\paragraph{Recorded target.}#1}
\begin{document}
% UnsolvedMath ID 1142; source catalogue code DYN-002
\setcounter{section}{1141}
\section{The Painlevé Conjecture}\label{1142}
{\small\sffamily UnsolvedMath ID 1142 · Dynamical Systems}
\subsection{Definitions and mathematical statement}

\paragraph{Shared notation.}$\mathbb N=\{1,2,\ldots\}$, and $\mathbb Z,\mathbb Q,\mathbb R,\mathbb C$ denote the integers, rationals, reals and complex numbers. Any other conventions are specified locally.


For a positive integer \(n\), consider positive masses
\(m_1,\ldots,m_n\) and positions \(q_i(t)\in\mathbb R^3\)
obeying Newtonian gravity, with the gravitational constant normalized to one:
\[
 \ddot q_i(t)=\sum_{j\ne i}m_j
      \frac{q_j(t)-q_i(t)}{\|q_j(t)-q_i(t)\|^3}.
\]
Initial positions are pairwise distinct and initial velocities are finite.
A collision singularity at a finite terminal time is one at which the
positions have finite limits with two of those limits coinciding.
A noncollision singularity is a finite endpoint of a maximal classical
solution that is not such a collision singularity; Newtonian examples
have bodies escaping to unbounded distances in finite time.
There are no collisions at earlier times of the classical solution.
\paragraph{Statement recorded in the supplied source.}
For each integer \(n\ge4\), do there exist positive masses and
collision-free initial conditions for which the Newtonian \(n\)-body
solution has a noncollision singularity at some finite time \(T>0\)?
The spatial formulation includes planar motions.
The historical existence question has been answered affirmatively;
the cited four-body construction settles its final threshold case. \sref{1142}{2}
\subsection{Short English statement}
Can Newtonian point masses develop a finite-time singularity without a collision, for every number of bodies at least four? The cited work resolves the historical four-body case.
\subsection{Sources}
\begin{description}
\item[\hypertarget{source:1142:1}{S1}] ULAM.AI and the UnsolvedMath Contributors, \emph{UnsolvedMath: A Curated Collection of Open Mathematics Problems}, record 1142 (DYN-002). CC BY 4.0. \url{https://huggingface.co/datasets/ulamai/UnsolvedMath}.
\item[\hypertarget{source:1142:2}{S2}] Jinxin Xue, \emph{Noncollision Singularities in a Planar Four-body Problem}, Acta Mathematica 224 (2020), 253--388, Section 1.1, Conjecture 2, and Section 1, Theorem 1; arXiv:1409.0048. \url{https://arxiv.org/abs/1409.0048}.
\end{description}
\label{end:1142}
\end{document}
